\section{Cluster algebras and upper cluster algebras}\label{section2}

Let $m$, $n$ be positive integers such that $m\geq n$. Denote $\mathcal{F}=\mathbb{Q}(x_1,\dots,x_m)$.
A \emph{seed}  $\Sigma=(\tilde{\bf x},\tilde{B})$ is a pair where $\tilde{\bf x}=\{x_1,\ldots,x_m\}$ is an $m$-tuple of elements of $\mathcal{F}$ that form a free generating set, $\tilde{B}$ is an $m\times n$ integer matrix such that the submatrix $B$  (called the \emph{principal part}) formed by the top $n$ rows is sign-skew-symmetric (that is, either $b_{ij}=b_{ji}=0$, or else~$b_{ji}$ and~$b_{ij} $ are of opposite sign; in particular, $b_{ii}=0$ for all~$i$).
The integer $n$ is called the \emph{rank} of the seed.

For any integer $a$, let $[a]_+:= \max(0,a)$.
 Given a seed $(\tilde{\bf x},\tilde{B})$ and a specif\/ied index $1\leq k \leq n$, we def\/ine \emph{mutation} of $(\tilde{\bf x},\tilde{B})$ at $k$, denoted $\mu_{k}(\tilde{\bf x},\tilde{B})$, to be a new seed $(\tilde{\bf x}',\tilde{B}')$, where
\begin{gather*}
 x_i' = \begin{cases}
\displaystyle  x'_k=x_k^{-1}\left(\prod\limits_{i=1}^{m}x_i^{[b_{ik}]_+}+\prod\limits_{i=1}^{m}x_i^{[-b_{ik}]_+}\right),
  & \mbox{if} \ \  i=k, \\
x_i, & \mbox{otherwise}, \end{cases}\\
b_{ij}' = \begin{cases} -b_{ij}, & \mbox{if} \ \ i=k \ \ \text{or} \ \ j=k,  \\
b_{ij} + \dfrac{|b_{ik}|b_{kj} + b_{ik}|b_{kj}|}{2}, & \mbox{otherwise.}
\end{cases}
\end{gather*}
If the principal part of $\tilde{B}'$ is also sign-skew-symmetric, we say that the mutation is well-def\/ined.
Note that a well-def\/ined mutation is an involution, that is mutating $(\tilde{\bf x}', \tilde{B}')$ at~$k$ will return to our original seed $(\tilde{\bf x}, \tilde{B})$.

Two seeds $\Sigma_1$ and $\Sigma_2$ are said to be \emph{mutation-equivalent} or in the same \emph{mutation class} if~$\Sigma_2$ can be obtained by a sequence of well-def\/ined mutations from~$\Sigma_1$. This is obviously an equivalence relation. A seed~$\Sigma$ is said to be \emph{totally mutable}  if every sequence of mutations from~$\Sigma$ consists of well-def\/ined ones. It is shown in \cite[Proposition~4.5]{FZ1} that a seed is totally mutable if $B$ is skew-symmetrizable, that is, if there exists a diagonal matrix $D$ with positive diagonal entries such that $DB$ is skew-symmetric.


To emphasize the dif\/ferent roles played by $x_i$ $(i\le n)$ and $x_i$ $(i>n)$, we also use $({\bf x},{\bf y},B)$ to denote the seed $(\tilde{\bf x}, \tilde{B})$, where ${\bf x}=\{x_1,\dots,x_n\}$, ${\bf y}=\{y_1,\dots,y_n\}$ where $y_j=\prod\limits_{i=n+1}^mx_i^{b_{ij}}$. We call~${\bf x}$ a~\emph{cluster}, ${\bf y}$ a~\emph{coefficient tuple}, $B$ the \emph{exchange matrix}, and the elements of a cluster \emph{cluster variables}. We denote
\begin{gather*}
\mathbb{ZP}=\mathbb{Z}\big[x_{n+1}^{\pm1},\dots,x_m^{\pm1}\big].
\end{gather*}

In the paper, we shall only study cluster algebras of geometric type, def\/ined as follows.

\begin{Definition}
Given a totally mutable seed $({\bf x}, {\bf y},B)$, the \emph{cluster algebra} $\mathcal{A}({\bf x},{\bf y},B)$ \emph{of geometric type} is the subring of $\mathcal{F}$ generated over $\mathbb{ZP}$ by
\begin{gather*}
\bigcup_{({\bf x}',{\bf y}',B')} {\bf x}',
 \end{gather*}
 where the union runs over all seeds $({\bf x}', {\bf y}',B')$ that are mutation-equivalent to $({\bf x},{\bf y},B)$. The seed $({\bf x},{\bf y},B)$ is called the \emph{initial seed} of $\mathcal{A}({\bf x},{\bf y},B)$. (Since $({\bf x},{\bf y},B)=(\tilde{\bf x},\tilde{B})$ in our notation, $\mathcal{A}({\bf x},{\bf y},B)$ is also denoted $\mathcal{A}(\tilde{\bf x},\tilde{B})$.)
\end{Definition}

It follows from the def\/inition that any seed in the same mutation class will generate the same cluster algebra up to isomorphism.

 For any $n \times n$ sign-skew-symmetric matrix $B$, we associate a (simple) directed graph $Q_B$ with vertices $1,\ldots,n$, such that for each pair $(i,j)$ with $b_{ij}>0$, there is  \emph{exactly one} arrow from vertex~$i$ to vertex~$j$.
(Note that even if $B$ is skew-symmetric, $Q_B$ is not the usual quiver associated to~$B$ which can have multiple edges.)


 We call $B$ (as well as the digraph $Q_B$ and the seed $\Sigma=({\bf x},{\bf y},B)$) \emph{acyclic} if there are no oriented cycles in $Q_B$.
We say that the cluster algebra $\mathcal{A}({\bf x},{\bf y},B)$ is \emph{acyclic} if there exists an acyclic seed; otherwise we say that the cluster algebra is \emph{non-acyclic}.

\begin{Definition}
Given a cluster algebra $\mathcal{A}$, the \emph{upper cluster algebra} $\mathcal{U}$ is def\/ined as
\begin{gather*}
 \mathcal{U} =
 \bigcap_{{\bf x}=\{x_1,\ldots,x_n\}}
  \mathbb{ZP}\big[x_1^{\pm 1},\ldots,x_n^{\pm 1}\big],
\end{gather*}
where ${\bf x}$ runs over all clusters of~$\mathcal{A}$.
\end{Definition}

Now we can give the following def\/inition of coprime when the cluster algebra is of geometric type (given in \cite[Lemma~3.1]{BFZ}).

\begin{Definition}
A seed  $(\tilde{\bf x},\tilde{B})$ is \emph{coprime} if no two columns of $\tilde{B}$ are proportional to each other with the proportionality coef\/f\/icient being a ratio of two odd integers.

A cluster algebra is \emph{totally coprime} if every seed is coprime.
\end{Definition}


In certain cases it is suf\/f\/icient to consider only the clusters of the initial seed and the seeds that are a single mutation away from it, rather than all the seeds in the entire mutation class.
For a cluster ${\bf x}$,  let $\mathcal{U}_{{\bf x}}$ be the intersection in $\mathbb{ZP}(x_1,\ldots,x_n)$ of the $n+1$ Laurent rings corresponding to ${\bf x}$ and its one-step mutations:
\begin{gather*}
\mathcal{U}_{\bf x} := \mathbb{ZP}\big[x_1^{\pm 1},\ldots, x_n^{\pm 1}\big] \cap \left(\bigcap_i \mathbb{ZP}\big[x_1^{\pm 1},\ldots,x_i'^{\pm 1},\ldots, x_n^{\pm 1}\big]\right).
\end{gather*}

\begin{Theorem}[\cite{BFZ,M}]\label{UB}
We have
 $\mathcal{A}\subseteq\mathcal{U}\subseteq\mathcal{U}_{\bf x}$. Moreover,
\begin{enumerate}\itemsep=0pt
\item[$(i)$] If $\mathcal{A}$ is acyclic, then $\mathcal{A}=\mathcal{U}$.

\item[$(ii)$] If $\mathcal{A}$ is totally coprime, then $\mathcal{U}=\mathcal{U}_{{\bf x}}$ for any seed $({\bf x},{\bf y},B)$. In particular, this holds when the matrix~$\tilde{B}$ has full rank.
\end{enumerate}
\end{Theorem}


